% SEGMENT OLP-0097-S01
% Part: first-order-logic
% Chapter: natural-deduction
% Section: soundness-identity

\documentclass[../../../include/open-logic-section]{subfiles}

\begin{document}

\olfileid{fol}{ntd}{sid}

\olsection{\usetoken{S}{identity} கொண்ட சரித்தன்மை}

\begin{prop}
$\eq$ க்கான விதிகளைக் கொண்ட இயல்பான வருவித்தல் சரித்தன்மையுடையது.
\end{prop}

\begin{proof}
$\eq[t][t]$ வடிவிலான எந்த !!{formula} உம் செல்லுபடியாகும்; ஏனெனில் ஒவ்வொரு
!!{structure}~$\Struct M$ க்கும் $\Sat{M}{\eq[t][t]}$. (சொல்~$t$ மூடியது,
அதாவது எந்த மாறியையும் கொண்டிருக்காது என்று எடுத்துக்கொள்கிறோம்; எனவே மாறி
ஒதுக்கீடுகள் பொருட்டல்ல.)

% SEGMENT OLP-0097-S02
!!a{derivation}[gen] கடைசி அனுமானம் \Elim{\eq} என்க; அதாவது !!{derivation}
பின்வரும் வடிவில் உள்ளது:
\begin{prooftree}
  \AxiomC{$\Gamma_1$}
  \RightLabel{$\delta_1$}
  \DeduceC{$\eq[t_1][t_2]$}
  \AxiomC{$\Gamma_2$}
  \RightLabel{$\delta_2$}
  \DeduceC{$!A(t_1)$}
  \RightLabel{\Elim{\eq}}
  \BinaryInfC{$!A(t_2)$}
\end{prooftree}
முன்கோள்கள்~$\eq[t_1][t_2]$, $!A(t_1)$ ஆகியவை முறையே !!{undischarged}
எடுகோள்கள்~$\Gamma_1$, $\Gamma_2$ இலிருந்து !!{derive} செய்யப்படுகின்றன.
$!A(t_2)$ ஆனது $\Gamma_1 \cup \Gamma_2$ இலிருந்து பின்தொடர்கிறது என்று காட்ட
விரும்புகிறோம். $\Sat{M}{\Gamma_1 \cup \Gamma_2}$ ஆகுமாறு
!!a{structure}~$\Struct{M}$ ஐக் கருதுக. தொகுத்தறிதல் கருதுகோளால்
$\Sat{M}{!A(t_1)}$ மற்றும் $\Sat{M}{\eq[t_1][t_2]}$. ஆகவே
$\Value{t_1}{M} = \Value{t_2}{M}$. $s$ ஏதேனும் ஒரு மாறி ஒதுக்கீடு என்றும்,
$m = \Value{t_1}{M} = \Value{t_2}{M}$ என்றும் கொள்க.
\oliflabeldef{fol:syn:ext:prop:ext-formulas}{\olref[fol][syn][ext]{prop:ext-formulas}}{சொல்
இடைநிறுத்தலுக்கும் மாறி ஒதுக்கீட்டுக்கும் இடையிலான முன்மொழிவு} படி,
$\Sat{M}{!A(t_1)}[s]$ என்பது, அப்போதும் அப்போது மட்டுமே,
$\Sat{M}{!A(x)}[\Subst{s}{m}{x}]$; அதுவும், அப்போதும் அப்போது மட்டுமே,
$\Sat{M}{!A(t_2)}[s]$. $\Sat{M}{!A(t_1)}$ என்பதால்,
$\Sat{M}{!A(t_2)}$.
\end{proof}

\end{document}
